\newcommand{\titulus}{\nomenFesti{S. Henrici.}
\dies{Die 13. Iulii.}}
\newcommand{\oratio}{\pars{Oratio.}

\noindent Deus, qui beátum Henrícum, grátiæ tuæ ubertáte prævéntum, e terréni cura regíminis ad supérna mirabíliter erexísti, eius nobis intercessióne largíre, ut inter mundánas varietátes puris ad te méntibus festinémus.

\pars{Pro pace in universo mundo.} \scriptura{Sir. 50, 25; 2 Esdr. 4, 20; \textbf{H416}}

\vspace{-4mm}

\antiphona{II D}{temporalia/ant-dapacemdomine.gtex}

\vfill

\noindent Deus, a quo sancta desidéria, recta consília et iusta sunt ópera: da servis tuis illam, quam mundus dare non potest, pacem; ut et corda nostra mandátis tuis dédita, et hóstium subláta formídine, témpora sint tua protectióne tranquílla.

\noindent Per Dóminum nostrum Iesum Christum, Fílium tuum, qui tecum vivit et regnat in unitáte Spíritus Sancti, Deus, per ómnia sǽcula sæculórum.

\noindent \Rbardot{} Amen.}
\newcommand{\invitatorium}{\pars{Invitatorium.}

\vspace{-4mm}

\antiphona{IV}{temporalia/inv-mirabileminsanctis.gtex}}
\newcommand{\hymnusmatutinum}{\pars{Hymnus}

\cuminitiali{VIII}{temporalia/hym-IesuRedemptor.gtex}}
\newcommand{\matversus}{\noindent \Vbardot{} Iustum dedúxit Dóminus per vias rectas.

\noindent \Rbardot{} Et osténdit illi regnum Dei.}
\newcommand{\lectioi}{\pars{Lectio I.} \scriptura{1 Reg. 18, 16b-40}

\noindent De libro primo Regum.

\noindent In diébus illis: Venit Achab in occúrsum Elíæ et, cum vidísset eum, ait: «Tune es, qui contúrbas Israel?».

\noindent Et ille ait: «Non turbávi Israel, sed tu et domus patris tui, qui dereliquístis mandáta Dómini, et secútus es Báalim. Verúmtamen nunc mitte et cóngrega ad me univérsum Israel in monte Carméli et prophétas Baal quadringéntos quinquagínta prophetásque Aseræ quadringéntos, qui cómedunt de mensa Iézabel».

\noindent Misit Achab ad omnes fílios Israel et congregávit prophétas in monte Carméli.

\noindent Accédens autem Elías ad omnem pópulum ait: «Usquequo claudicátis in duas partes? Si Dóminus est Deus, sequímini eum; si autem Baal, sequímini illum».

\noindent Et non respóndit ei pópulus verbum.

\noindent Et ait rursus Elías ad pópulum: «Ego remánsi prophéta Dómini solus; prophétæ autem Baal quadringénti et quinquagínta viri sunt. Dentur nobis duo boves, et illi éligant sibi bovem unum et in frusta cædéntes ponant super ligna; ignem autem non suppónant. Et ego fáciam bovem álterum et impónam super ligna; ignémque non suppónam. Invocáte nomen dei vestri, et ego invocábo nomen Dómini; et Deus, qui exaudíerit per ignem, ipse est Deus!». Respóndens omnis pópulus ait: «Optima proposítio».

\noindent Dixit ergo Elías prophétis Baal: «Elígite vobis bovem unum et fácite primi, quia vos plures estis; et invocáte nomen dei vestri ignémque non supponátis». Qui cum tulíssent bovem, quem déderat eis, fecérunt et invocábant nomen Baal de mane usque ad merídiem dicéntes: «Baal, exáudi nos!». Et non erat vox, nec qui respondéret.

\noindent Saliebántque in circúitu altáris, quod fécerant. Cumque esset iam merídies, illudébat eis Elías dicens: «Clamáte voce maióre; deus enim est et fórsitan occupátus est aut secéssit aut in itínere aut certe dormit, ut excitétur».

\noindent Clamábant ergo voce magna et incidébant se iuxta ritum suum cultris et lancéolis, donec perfunderéntur sánguine.}
\newcommand{\responsoriumi}{\pars{Responsorium 1.} \scriptura{\Rbardot{} 1 Reg. 7, 3 \Vbardot{} ibidem; \textbf{H397}}

\vspace{-5mm}

\responsorium{III}{temporalia/resp-praeparatecordavestra-sinedox.gtex}{}}
\newcommand{\lectioii}{\pars{Lectio II.} \scriptura{1 Reg. 18, 29-40}

\noindent Postquam autem, tránsiit merídies et, illis prophetántibus, vénerat tempus, quo sacrifícium offérri solet, nec audiebátur vox, neque áliquis respondébat nec attendébat orántes, dixit Elías omni pópulo: «Veníte ad me».

\noindent Et, accedénte ad se pópulo, curávit altáre Dómini, quod destrúctum fúerat; et tulit duódecim lápides iuxta númerum tríbuum filiórum Iacob, ad quem factus est sermo Dómini dicens: «Israel erit nomen tuum».

\noindent Et ædificávit lapídibus altáre in nómine Dómini fecítque aquædúctum quasi pro duóbus satis in circúitu altáris et compósuit ligna divisítque per membra bovem et pósuit super ligna et ait: «Impléte quáttuor hýdrias aqua et fúndite super holocáustum et super ligna».

\noindent Rursúmque dixit: «Etiam secúndo hoc fácite». Qui cum fecíssent et secúndo, ait: «Etiam tértio idípsum fácite». Fecerúntque et tértio, et currébant aquæ circum altáre, et fossa aquædúctus repléta est.

\noindent Cumque iam tempus esset, ut offerrétur sacrifícium, accédens Elías prophéta ait: «Dómine, Deus Abraham, Isaac et Israel, hódie osténde quia tu es Deus in Israel et ego servus tuus et iuxta præcéptum tuum feci ómnia hæc. Exáudi me, Dómine, exáudi me, ut discat pópulus iste quia tu, Dómine, es Deus et tu convertísti cor eórum íterum!».

\noindent Cécidit autem ignis Dómini et vorávit holocáustum et ligna et lápides, púlverem quoque et aquam, quæ erat in aquædúctu lambens. Quod cum vidísset omnis pópulus, cécidit in fáciem suam et ait: «Dóminus ipse est Deus, Dóminus ipse est Deus!».

\noindent Dixítque Elías ad eos: «Apprehéndite prophétas Baal, et ne unus quidem effúgiat ex eis!». Quos cum comprehendíssent, duxit eos Elías ad torréntem Cison et interfécit eos ibi.}
\newcommand{\responsoriumii}{\pars{Responsorium 2.} \scriptura{\Vbardot{} Ex. 34, 29; \textbf{H160}}

\vspace{-5mm}

\responsorium{VIII}{temporalia/resp-splendidafactaestfacies-CROCHU.gtex}{}}
\newcommand{\lectioiii}{\pars{Lectio III.} \scriptura{MGH, Scriptores 4, 792-799}

\noindent E Vita antíqua sancti Henríci.

\noindent Unctus in regem beatíssimus Dei fámulus, temporális regni non conténtus angústiis, pro adipiscénda immortalitátis coróna, summo Regi, cui servíre regnáre est, militáre dispósuit. Summam étenim diligéntiam in amplificándo cultu religiónis adhíbuit; ecclésias ditáre possessiónibus et imménsis ornátibus augére cœpit. Episcopátum Bambergénsem ex íntegro in suo dómate fundávit, eundémque princípibus Apostolórum Petro et Paulo et pretiosíssimo mártyri Geórgio attitulátum, speciáli iure sanctæ Románæ Ecclésiæ contrádidit, ut et primæ sedi débitum divínitus impénderet honórem et suam plantatiónem tanto patrocínio fírmius communíret.

\noindent Ut autem cunctis liquídius enitéscat, qua vigilántia vir beatíssimus novéllæ suæ Ecclésiæ bona pacis et tranquillitátis étiam per succedéntia témpora províderit, áliquam hic insérimus confirmatiónis epístolam.

\noindent «Henrícus, divína præordinánte cleméntia rex, ómnibus Ecclésiæ fíliis, tam futúris quam præséntibus. Salubérrimis sacri elóquii institutiónibus erudímur et præmonémur, ut temporália relinquéntes bona, et terréna postponéntes cómmoda, ætérna et sine fine mansúra in cælis studeámus adipísci consistória. Glória enim præsens fugitíva est et vana, dum possidétur, nisi in ea áliquid de cælésti æternitáte cogitétur. Sed Dei miserátio humáno géneri útile provídit remédium, quando partem cæléstis pátriæ terrénæ substántiæ fecit esse prétium.

\noindent Huius ergo nos cleméntiæ non immémores, nec ignorántes nos gratuíto divínæ miseratiónis respéctu regáli dignitáte sublimátos, cóngruum esse dúcimus, non solum ecclésias ab antecessóribus nostris constrúctas ampliáre, sed ad maiórem Dei glóriam novas ædificáre eásque devotiónis nostræ donis gratíssimis exaltáre. Quaprópter domínicis non surdum audítum præbéntes præcéptis et deíficis obtemperándo intendéntes suasiónibus, thesáuros divínæ largitátis munificéntia nobis collátos, in cælo desiderámus repónere, ubi neque fures effódiant neque furéntur, neque ærúgo vel tínea demoliátur, ubi et, dum ómnia nunc congésta recólimus, cor nostrum desidério et amóre sǽpius versétur.

\noindent Proínde patére vólumus ómnium fidélium universitáti, quod quemdam patérnæ hereditátis nostræ locum, Babenberch dictum, in sedem et culmen episcopátus provéximus, ut ínibi nostrum parentúmque nostrórum célebre habeátur memoriále, et iugis pro ómnibus orthodóxis mactétur hóstia salutáris».}
\newcommand{\responsoriumiii}{\pars{Responsorium 3.} \scriptura{\Rbardot{} Mt. 25, 21 \Vbardot{} ibid., 20; \textbf{H377}}

\vspace{-5mm}

\responsorium{VII}{temporalia/resp-eugeservebone-CROCHU-cumdox.gtex}{}}
\newcommand{\hymnuslaudes}{\pars{Hymnus}

\cuminitiali{VIII}{temporalia/hym-IesuCorona.gtex}}
\newcommand{\lectiobrevis}{\pars{Lectio Brevis.} \scriptura{Rom. 12, 1-2 }

\noindent Obsecro vos, fratres, per misericórdiam Dei, ut exhibeátis córpora vestra hóstiam vivéntem, sanctam, Deo placéntem, rationábile obséquium vestrum; et nolíte conformári huic sǽculo, sed transformámini renovatióne mentis, ut probétis quid sit volúntas Dei, quid bonum et bene placens et perféctum.}
\newcommand{\responsoriumbreve}{\pars{Responsorium breve.} \scriptura{Sir. 45, 9}

\antiphona{VI}{temporalia/resp-amaviteum.gtex}}
\newcommand{\preces}{\noindent Christum Deum sanctum, fratres, exaltémus,~\gredagger{} orántes ut serviámus illi in sanctitáte et iustítia coram ipso ómnibus diébus nostris,~\grestar{} et acclamémus:

\Rbardot{} Tu solus sanctus, Dómine.

\noindent Qui tentári voluísti per ómnia pro similitúdine nostra absque peccáto,~\grestar{} miserére nostri, Dómine Iesu.

\Rbardot{} Tu solus sanctus, Dómine.

\noindent Qui nos omnes ad perfectiónem caritátis vocásti,~\grestar{} sanctífica nos, Dómine Iesu.

\Rbardot{} Tu solus sanctus, Dómine.

\noindent Qui nos iussísti esse salem terræ et lucem mundi,~\grestar{} illúmina nos, Dómine Iesu.

\Rbardot{} Tu solus sanctus, Dómine.

\noindent Qui voluísti ministráre,~\gredagger{} non ministrári,~\grestar{} fac nos tibi et frátribus humíliter servíre, Dómine Iesu.

\Rbardot{} Tu solus sanctus, Dómine.

\noindent Tu, splendor glóriæ Patris et figúra substántiæ eius,~\grestar{} fac ut in glória vultum tuum respiciámus, Dómine Iesu.

\Rbardot{} Tu solus sanctus, Dómine.}
\newcommand{\benedictus}{\pars{Canticum Zachariæ.} \scriptura{Rom. 8, 28; \textbf{H333}}

\vspace{-6mm}

\antiphona{VII d}{temporalia/ant-scimusquoniam.gtex}

\vspace{-3mm}

\scriptura{Lc. 1, 68-79}

\vspace{-3mm}

\cantusSineNeumas
\initiumpsalmi{temporalia/benedictus-initium-vii-d3-auto.gtex}

\vspace{-1.5mm}

\input{temporalia/benedictus-vii-d3.tex} \Abardot{}}
\newcommand{\benedicamuslaudes}{\cuminitiali{}{temporalia/benedicamus-memoria-laudes.gtex}}
\include{hebdomadaxv}
\include{feriaii}
